\section{The local two-flight law}
\label{sec:local}
Put $o=1-b$. With longitudinal displacement taken into the corresponding
obstacles, one flight has length
\[
 \ell_{bo}(u,w)=
 \sqrt{\{g+\psi_b(u)+\psi_o(w)\}^2+(u-w)^2}.
\]
The broken action is
$\ell_{bo}(u,w)+\ell_{ob}(w,v)$. Its $w$ derivative vanishes at
$(u,w,v)=(0,0,0)$, with second $w$ derivative $2c_o/g>0$.
The implicit-function theorem gives a unique stationary point
$w=w_b(u,v)$ on a fixed endpoint box. The segment clearance and
reflection law ensure that it is the actual local three-impact path.
Let
\[
 W_b(u,v)=\ell_{bo}(u,w_b)+\ell_{ob}(w_b,v),\qquad
 E_b=W_b-2g.
\]

\begin{lemma}[Quadratic action and physical density]\label{lem:flux}
Writing $y=(u,v)^t$, one has
\begin{equation}\label{eq:quad}
 w_b=\frac{u+v}{2c_o}+O(|y|^2),\qquad
 E_b(y)=E_{b,0}(y)+O(|y|^3),\qquad
 E_{b,0}(y)=\tfrac12 y^tH_by,
\end{equation}
where
\begin{equation}\label{eq:H}
 H_b=\frac1g
 \begin{pmatrix}c_b-(2c_o)^{-1}&-(2c_o)^{-1}\\
                 -(2c_o)^{-1}&c_b-(2c_o)^{-1}\end{pmatrix},
 \quad a_b=\sqrt{\det H_b}=\frac1g\sqrt{\frac{c_bz}{c_o}}.
\end{equation}
The normalized twist
\begin{equation}\label{eq:twist}
 b_b(y)=\frac{-W_{b,uv}(y)}{d_b^0},\qquad
 d_b^0=\frac1{2gc_o}=\frac{a_b}{\sinh(2\gamma)}
\end{equation}
is positive on a smaller box and equals one at the origin. On a
sufficiently small fixed offset collar,
\begin{align}
 G_b(d)&=\frac{a_b}{\pi d^2}
             \int(d-E_b(y))_+b_b(y)\dd y,\label{eq:Gintegral}\\
 J_b(d)&=\frac{a_b}{\pi d^2}
             \int(u+v)(d-E_b(y))_+b_b(y)\dd y.\label{eq:Jintegral}
\end{align}
Both functions extend smoothly to zero, with $G_b(0)=1$ and
$J_b(0)=0$. The constructions and every fixed derivative are uniform
on compact local smooth families with the necessary higher bounds.
\end{lemma}
\begin{proof}
The quadratic broken action is
\[
 \frac{c_bu^2+2c_ow^2+c_bv^2-2w(u+v)}{2g}.
\]
Eliminating $w$ proves \eqref{eq:quad}--\eqref{eq:H}.
Both eigenvalues of $H_b$ are positive, since $c_bc_o>1$.
The mixed derivative at zero gives \eqref{eq:twist}.
In the coordinates of a collision section, Liouville volume has density
$\dd u\dd p\dd r$, where $p$ is momentum conjugate to $u$ and $r$
is time before the first impact. The generating-function identity
$p=-W_{b,u}$ changes its density to
$(-W_{b,uv})\dd u\dd v\dd r$ on this dispersing branch.
The time constraint is $0<r<d-E_b(u,v)$.

Choose the collar smaller than a common positive lower bound on free
flight lengths. Then the initial residual interval does not cross a
preceding collision, and the time remaining after the third collision
cannot reach a fourth one. The strict local minimum of $E_b$ places
its entire active sublevel inside the fixed facing patches when the
collar is small. Every parametrized path is therefore a physical first-hit
path in the specified event; conversely every such near-onset path is
in this chart. Division by $2\pi A$ and integration in $r$ give the raw
versions of \eqref{eq:Gintegral} and \eqref{eq:Jintegral}.
Multiplication by $2A\sinh(2\gamma)/d^2$ proves the displayed formulas.
The same calculation holds with the bounded mark $u+v$, without
changing the measure or conditioning on success.

For completeness, let $y=\Phi_b(x)$ be a local smooth Morse chart,
with $E_b(\Phi_b(x))=|x|^2/2$ and $D\Phi_b(0)=H_b^{-1/2}$.
After $x=\sqrt d\,X$, division of either integral by $d^2$ leaves
an integral on the fixed disk $|X|^2<2$ with weight $1-|X|^2/2$.
For any smooth amplitude its Taylor expansion in $h=\sqrt d$ has only
even integrated powers: the integral at $-h$ equals that at $h$ by
$X\mapsto-X$. Taylor's formula, to arbitrarily high fixed order,
therefore gives a smooth function of $d\ge0$, including parameter
derivatives. For analytic graphs the same argument applies to a
convergent local series. The constant amplitude has integral
\[
 I_0(d):=\int(d-E_{b,0})_+\dd y=\frac{\pi d^2}{a_b}.
\]
This gives $G_b(0)=1$. The marked amplitude vanishes at the origin,
so $J_b(0)=0$ and $J_b(d)=O(d)$.
\end{proof}

The first part of Theorem~\ref{thm:main} follows immediately. Notice
that area is extracted from a raw exact germ before the normalization
in \eqref{eq:coefficients}; no unknown area is inserted into a finite
observation.
